\section{AI That Reads and Writes}
\begin{frame}{}
    \LARGE Part 5: \textbf{AI That Reads and Writes}
\end{frame}

\begin{frame}{Language is Hard for Computers}
    \textit{We understand each other so easily that we forget how tricky language is.}

    \vspace{0.8em}
    \begin{itemize}
        \setlength\itemsep{0.9em}
        \item \textbf{One word, many meanings}\\
        ``I went to the \textbf{bank}.'' --- to get money, or to sit by the river?
        \item \textbf{One sentence, two readings}\\
        ``I saw the man with the telescope.'' --- who has the telescope?
        \item \textbf{Meaning beyond the words}\\
        ``Great, another Monday.'' --- is the speaker happy?
    \end{itemize}

    \vspace{0.8em}
    \takeaway{Teaching computers to work with human language is called
    \textbf{Natural Language Processing (NLP)}.}
\end{frame}

\begin{frame}{What Language AI Can Do}
    \centering
    \renewcommand{\arraystretch}{1.5}
    \begin{tabular}{@{} >{\raggedright\arraybackslash}p{0.24\textwidth}
                        >{\raggedright\arraybackslash}p{0.30\textwidth}
                        >{\raggedright\arraybackslash}p{0.36\textwidth} @{}}
        \toprule
        \textbf{Task} & \textbf{What It Does} & \textbf{Example} \\
        \midrule
        \textbf{Translation} & Language to language & English $\rightarrow$ Arabic \\
        \textbf{Sentiment Analysis} & Detect mood & ``This movie was amazing!'' $\rightarrow$ Positive \\
        \textbf{Summarisation} & Shorten a long text & A 20-page report $\rightarrow$ 5 bullet points \\
        \textbf{Question Answering} & Find or write an answer & ``When does the library close?'' \\
        \textbf{Speech Recognition} & Turn voice into text & Voice messages, subtitles \\
        \bottomrule
    \end{tabular}
\end{frame}

\begin{frame}{Large Language Models: Guess the Next Word}
    \centering
    {\Large ``The capital of France is \underline{\hspace{2cm}}''}

    \vspace{0.8em}
    \begin{tikzpicture}[font=\small]
        \foreach \w/\len/\pct [count=\i] in {Paris/7.0/very likely, a/0.9/unlikely, located/0.6/unlikely, beautiful/0.4/unlikely} {
            \node[anchor=east] at (0,-\i*0.62) {\w};
            \fill[KAUSTblue!70] (0.15,-\i*0.62-0.18) rectangle (0.15+\len,-\i*0.62+0.18);
            \node[anchor=west, font=\scriptsize, gray] at (0.25+\len,-\i*0.62) {\pct};
        }
    \end{tikzpicture}

    \vspace{0.6em}
    \raggedright
    \begin{itemize}
        \setlength\itemsep{0.4em}
        \item A \textbf{Large Language Model (LLM)} is a huge neural network trained on one simple game: \textbf{predict the next word}.
        \item It plays this game on an enormous amount of text --- books, websites, articles.
        \item To get good at it, the model has to pick up grammar, facts and even some reasoning.
    \end{itemize}
\end{frame}

\begin{frame}{From Guessing Words to Having a Conversation}
    \begin{columns}[c]
        \begin{column}{0.5\textwidth}
            \centering
            \includegraphics[width=\linewidth,height=0.75\textheight,keepaspectratio]{images/foundation-models/chatgpt_example_conversation.png}
        \end{column}
        \begin{column}{0.47\textwidth}
            \begin{itemize}
                \setlength\itemsep{0.8em}
                \item You type a message.
                \item The model predicts a good next word, then the next, then the next\ldots
                \item One word at a time, a full answer appears.
                \item ChatGPT, Gemini, Claude and others all work this way.
            \end{itemize}
        \end{column}
    \end{columns}
\end{frame}

\begin{frame}{What Chat Assistants Are Good At}
    \begin{columns}[T]
        \begin{column}{0.48\textwidth}
            \begin{block}{Works well}
                \begin{itemize}
                    \setlength\itemsep{0.35em}
                    \item Drafting emails, letters and reports
                    \item Summarising long documents
                    \item Explaining a topic at your level
                    \item Translating and fixing grammar
                    \item Brainstorming ideas
                    \item Helping to write computer code
                \end{itemize}
            \end{block}
        \end{column}
        \begin{column}{0.48\textwidth}
            \begin{alertblock}{Be careful}
                \begin{itemize}
                    \setlength\itemsep{0.35em}
                    \item Exact facts, numbers and quotes
                    \item Very recent events
                    \item Medical, legal or financial decisions
                    \item Anything involving private information
                \end{itemize}
            \end{alertblock}
        \end{column}
    \end{columns}

    \vspace{0.8em}
    \centering
    \textit{Think of it as a fast, well-read assistant --- whose work you still need to check.}
\end{frame}

\begin{frame}{Asking Well: Prompting}
    \textit{The text you give the model is called a \textbf{prompt}. Better prompts give better answers.}

    \vspace{0.6em}
    \begin{columns}[T]
        \begin{column}{0.4\textwidth}
            \begin{enumerate}
                \setlength\itemsep{0.5em}
                \item Say \textbf{who it is for}
                \item Be \textbf{specific} about the task
                \item Give \textbf{context}
                \item Ask for a \textbf{format}
                \item \textbf{Improve} it step by step
            \end{enumerate}
        \end{column}
        \begin{column}{0.57\textwidth}
            \begin{tcolorbox}[colback=red!5, colframe=red!50, title=\textbf{Vague}]
                \small ``Write about exercise.''
            \end{tcolorbox}
            \begin{tcolorbox}[colback=blue!5, colframe=blue!50, title=\textbf{Clear}]
                \small ``Write a friendly 5-point list explaining to office workers why a
                20-minute daily walk is good for them. Keep each point to one sentence.''
            \end{tcolorbox}
        \end{column}
    \end{columns}
\end{frame}
